\section{The Faith of the Church}\label{sec:faith}

The Church confesses the angels in her Creed, and across sixteen centuries
her councils, popes, and catechisms have unfolded that confession: against
the dualists who made the devil a second principle beside God, against the
Origenists who made angels and souls one family of fallen minds, and against
the modern reduction of the angels to figures of speech. Each act is given
below in order of date, with its locus and its teaching.
Appendix~\ref{app:chronology} gathers the same acts, with the rulings that
govern the cult of the angels, in a single table.

\sectionguard
\subsection{The Creed: Maker of Things Invisible}

The Creed of Nicaea (325) and the creed received as that of the Council of
Constantinople (381) confess the one God as maker of all things visible and
invisible. The Constantinopolitan text, in the Latin that the Roman Church
uses, reads \latin{factorem caeli et terrae, visibilium omnium et
invisibilium} (DS~150). Under the ``things invisible'' the Church confesses
the angels. The whole unseen creation stands under the one Creator, and no
spirit is a second principle beside God or an emanation of his substance.
Every later act on the angels unfolds this article, and Paul VI's
\emph{Credo} of 1968 names what it has always meant: the things invisible
are ``the pure spirits which are also called angels''.

\sectionguard
\subsection{The Patristic Acts: Leo, Constantinople, Braga}

\paragraph{Leo I to Turibius (447).} In the letter \latin{Quam laudabiliter}
to Turibius, bishop of Astorga (21 July 447), Leo the Great answered a list
of Priscillianist errors. The sixth held that the devil was never good,
that his nature is not God's work, and that he ``emerged from chaos and
darkness'' as the principle and substance of evil. Leo replies that the true
faith confesses the substance of every creature, spiritual or corporeal, to
be good, and that evil has no nature, ``because God, who is the founder of
all things, made nothing that is not good'' (DS~286). He then states the
doctrine that the later councils take up:
\begin{quote}
\latin{Unde et diabolus bonus esset, si in eo quod factus est permaneret.
Sed quia naturali excellentia male usus est `et in veritate non stetit'
(Jo 8, 44), non in contrariam transiit substantiam, sed a summo bono, cui
debuit adhaerere, descivit.} (DS~286)
\end{quote}
The devil would be good if he had remained in what he was made. Because he
misused his natural excellence, he did not pass into a contrary substance;
he withdrew from the highest good to which he ought to have held fast. Leo's
teaching passed into the canons of Braga and, through them, into the
Lateran confession.

\paragraph{Constantinople (543).} The edict of the Emperor Justinian to
Menas, patriarch of Constantinople, was published in a synod of
Constantinople in 543 with nine anathematisms against Origen and his
followers (DS~403--411). Five touch the angels. The first condemns the teaching that human souls
pre-existed as ``minds and holy powers'' (\latin{mentes \dots\ et sanctae
virtutes}) who grew sated with the contemplation of God, cooled in charity,
and were sent into bodies as punishment (DS~403). The error made angels and
souls one class of spirits differing only in the degree of their fall; the
canon keeps them distinct in origin and in destiny. The fourth condemns the
claim that the Word became like all the heavenly orders, a cherub among the
cherubim and a seraph among the seraphim (DS~406): the Word took human
nature, not the nature of any angel. The sixth condemns the teaching that
the heavens, sun, moon, stars, and the waters above the heavens are
ensouled rational powers (DS~408). The seventh condemns the teaching that
Christ will be crucified in the world to come for demons as he was for men
(DS~409). The ninth condemns the teaching that the punishment of demons and
of impious men is temporary and will end in their restoration
(\latin{restitutionem et redintegrationem}, DS~411). These anathemas are
the earliest synodal witness that the punishment of the demons is without
end.

\paragraph{Braga I (561).} The first Council of Braga, a provincial synod of
the Suevic kingdom, issued a series of anathematisms, chiefly against the
Priscillianists. Four concern the angels. It condemns anyone who believes
that human souls or angels came from the substance of God (DS~455), or that
human souls first sinned in a heavenly dwelling and were cast down into
bodies (DS~456). The seventh canon repeats the error that Leo had answered
and anathematizes it:
\begin{quote}
\latin{Si quis dicit, diabolum non fuisse prius (bonum) angelum a Deo
factum, nec Dei opificium fuisse naturam eius, sed dicit eum ex chao et
tenebris emersisse nec aliquem sui habere auctorem, sed ipsum esse
principium atque substantiam mali \dots, an. s.} (DS~457)
\end{quote}
The eighth condemns the belief that the devil made some creatures in the
world and that he produces thunder, lightning, storms, and droughts by his
own authority (DS~458). The twelfth and thirteenth condemn the claims that
the forming of the human body in the womb is the work of demons and that
the creation of all flesh is the work of wicked angels (DS~462--463). The
demons can work in the world, as Scripture shows, but only as creatures
under God's permission; they create nothing and command nothing by right.
Against a dualist system the synod thus stated what the universal faith
already held and what Lateran~IV would define.

\paragraph{The early professions.} The \emph{Statuta Ecclesiae Antiqua}, a
collection from southern Gaul of the later fifth century, directs that a
bishop-elect be examined on whether he holds ``that the devil became evil
not by his condition but by his will'' (\latin{si diabolus non per
condicionem, sed per arbitrium factus sit malus}, DS~325). Innocent~III
put the same sentence into the profession of faith prescribed for Durand of
Osca and the Waldensians who returned to communion (\latin{Eius exemplo},
18~December 1208): \latin{Diabolum non per condicionem, sed per arbitrium
malum factum esse credimus} (DS~797). By the eve of the Fourth Lateran
Council, then, the goodness of the devil's created nature and the voluntary
origin of his evil were part of the confession required of a bishop and of
a reconciled heretic.

\sectionguard
\subsection{The Fourth Lateran Council (1215)}

The Fourth Lateran Council, convened by Innocent~III and held in November
1215, opened its constitutions with the profession \latin{Firmiter
credimus}. It was framed against the Cathars, who taught two principles and
ascribed the visible world to an evil god. The chapter confesses the one God
as
\begin{quote}
\latin{creator omnium visibilium et invisibilium, spiritualium et
corporalium: qui sua omnipotenti virtute simul ab initio temporis utramque
de nihilo condidit creaturam, spiritualem et corporalem, angelicam videlicet
et mundanam: ac deinde humanam, quasi communem ex spiritu et corpore
constitutam. Diabolus enim et alii daemones a Deo quidem natura creati sunt
boni, sed ipsi per se facti sunt mali. Homo vero diaboli suggestione
peccavit.} (DS~800)
\end{quote}
God, by his almighty power, at the beginning of time made from nothing both
the spiritual and the corporeal creature, the angelic and the earthly, and
afterwards the human creature, which shares in both, being made of spirit
and body. The devil and the other demons were created good in nature by
God; they became evil of themselves. Man sinned at the devil's suggestion.
The same chapter ends with the judgment, in which the reprobate will receive
``perpetual punishment with the devil'' (\latin{illi cum diabolo poenam
perpetuam}, DS~801).

In this solemn profession of an ecumenical council the Church defines five
truths concerning the angels. First, there are spiritual creatures distinct from the corporeal world,
and they are called angelic. Second, they were made by God from nothing,
and not from his substance or from any pre-existing stuff. Third, their
creation belongs to the beginning of time: angels are not eternal, and
their existence has a beginning. Fourth, the devil and the other demons are
creatures good by nature, whose evil is their own act and not their
nature. Fifth, their punishment, which the reprobate share, is perpetual.

The words \latin{simul ab initio temporis} confess that both creatures, the
spiritual and the corporeal, began with time. Aquinas reads them also of
the angels' creation together with the visible world and holds this the
more probable view, while the contrary opinion of many Greek Fathers, that
the angels were made before it, ``is not to be deemed erroneous'' (\ST{I
q.~61 a.~3}; Section~\ref{sec:creation}).

\sectionguard
\subsection{Florence and Trent}

\paragraph{Florence (1442).} The bull of union with the Copts,
\latin{Cantate Domino} (4 February 1442; 1441 in the Florentine reckoning),
confesses the one God as creator ``of all things visible and invisible'',
who ``made all creatures, both spiritual and corporeal: good, indeed,
because made by the highest good, but mutable, because they were made from
nothing'' (\latin{universas, tam spiritales quam corporales, condidit
creaturas: bonas quidem, quia a summo bono factae sunt, sed mutabiles, quia
de nihilo factae sunt}); and it
asserts that there is no nature of evil (\latin{nullamque mali asserit esse
naturam}, DS~1333). It then condemns the Manichaean doctrine of two first
principles, one of visible and the other of invisible things (DS~1336). The
decree adds one word to the Lateran confession, and that word bears on the
angels: their created goodness is mutable. A spirit made from nothing can
turn from the good, and the fall of some angels is thus possible without
any defect in their Maker. The council confirms what Lateran~IV had
defined.

\paragraph{Trent (1546).} The decree on original sin (Session~V, 17 June
1546) says that Adam, by his transgression, incurred ``captivity under the
power of him who thenceforth `had the empire of death' (Hebr 2,14), that
is, of the devil'' (\latin{captivitatem sub eius potestate, `qui mortis'
deinde `habuit imperium' \dots, hoc est diaboli}, DS~1511). In defining
original sin the council teaches the real power of the devil over fallen
man, a power held not by right but by the captivity that sin brought about.
Section~\ref{sec:assaults} returns to that power.

\sectionguard
\subsection{The First Vatican Council (1870)}

The dogmatic constitution \latin{Dei Filius} (24 April 1870) repeated the
Lateran text in its first chapter, adding that God created ``by his most
free counsel'' and ``not to increase his beatitude \dots\ but to manifest
his perfection'':
\begin{quote}
\latin{liberrimo consilio simul ab initio temporis utramque de nihilo
condidit creaturam, spiritualem et corporalem, angelicam videlicet et
mundanam, ac deinde humanam quasi communem ex spiritu et corpore
constitutam.} (DS~3002)
\end{quote}
The fifth canon on God the creator anathematizes whoever does not confess
that the world and all things in it, ``both spiritual and material'', were
brought forth by God from nothing ``according to their whole substance''
(\latin{et spirituales et materiales secundum totam suam substantiam a Deo
ex nihilo esse productas}, DS~3025). The canon was directed against
pantheism and against the theory that creation is necessary. It defines,
for the angels as for every creature, that nothing in them pre-exists their
creation. Vatican~I thus renewed the Lateran definition in the same words
and added the defined freedom of the act by which God made them.

\sectionguard
\subsection{\emph{Humani generis} (1950)}

In the encyclical \latin{Humani generis} (12 August 1950), Pius~XII listed
among the opinions that threatened the foundations of Catholic doctrine
that ``some also question whether angels are personal beings, and whether
matter and spirit differ essentially'' (n.~26; \latin{Quaestio etiam a
nonnullis agitur num Angeli creaturae personales sint}, DS~3891). Against writers who read the angels as personifications of
divine action or figures of religious language, Pius~XII closes the
question: the angels are personal beings. On this encyclical the Catechism
teaches that the angels are ``personal and immortal creatures'' (CCC~330,
citing DS~3891), endowed with intellect and will.

\sectionguard
\subsection{Paul VI: The Credo and the Audience of 1972}

\paragraph{The Credo of the People of God (1968).} At the close of the Year
of Faith, on 30 June 1968, Paul~VI pronounced a solemn profession of faith
that repeats in substance the creed of Nicaea, with the developments the
time required (\emph{Credo} 3). Its first article confesses God as creator ``of
things invisible such as the pure spirits which are also called angels''
(\emph{Credo} 8), and the Latin text cites DS~3002 at that point. Near its
end the \emph{Credo} places the saints in heaven ``associated with the holy
angels in the divine rule exercised by Christ in glory'' (\emph{Credo} 29).
The \emph{Credo} thus names what the Church has always understood by the
Creed's ``things invisible''.

\paragraph{The audience ``Deliver us from evil'' (1972).} At the general
audience of 15 November 1972, Paul~VI said that one of the greatest needs of
the Church was defence against ``that evil which we call the Devil'' (\latin{la
difesa da quel male, che chiamiamo il Demonio}). He described evil as ``no
longer only a deficiency, but an efficiency, a living being, spiritual,
perverted and perverting'' (\latin{un essere vivo, spirituale, pervertito e
pervertitore}), and added: ``Terrible reality. Mysterious and fearful.'' He
then drew the line: whoever refuses to recognize this reality as existing,
or makes of it a principle in its own right, ``goes outside the framework
of biblical and ecclesiastical teaching'' (\latin{Esce dal quadro
dell'insegnamento biblico ed ecclesiastico}). He cited DS~800 and added a
caution that guards the doctrine from excess: ``it is not said that every
sin is directly due to diabolical action'' (\latin{Non è detto che ogni
peccato sia direttamente dovuto ad azione diabolica}). The Pope's teaching
rests on the Lateran text, which he cites: to deny the devil's existence is
to depart from the Church's teaching.

\sectionguard
\subsection{\emph{Christian Faith and Demonology} (1975)}

On 26 June 1975 \emph{L'Osservatore Romano} published a study entitled
\emph{Fede cristiana e demonologia}, prepared by an expert at the request
of the Congregation for the Doctrine of the Faith, which in its prefatory
note strongly recommended the study as a sure basis for reaffirming the
doctrine of the magisterium.

The study reviews the witness of Scripture, the Fathers, and the councils,
and shows how the Church holds the existence of the demons. Lateran~IV,
framed against dualism, confesses that the devil and the other demons were
created good and became evil of themselves. That chapter is a single
confession, and ``each principal point of it has equal dogmatic value''
(\latin{ciascun punto principale di esso ha egualmente valore dogmatico}): a
text that teaches how the demons became evil teaches that they exist. In
the Church's constant faith and practice, above all in her liturgy and her
life of prayer, ``the existence of the demonic world is revealed as a
dogmatic datum'' (\latin{l'esistenza del mondo demoniaco si rivela come un
dato dogmatico}).

\sectionguard
\subsection{The Catecheses of John Paul II (1986)}

In the summer of 1986, within his catecheses on the Creed, John Paul~II
gave six general audiences on the angels and the demons (9, 23, and 30
July; 6, 13, and 20 August), the fullest papal catechesis on the angels.

The first audience (9 July) set the angels within the Creed: the truth
about them is ``collateral'' to the central revelation of God in Christ, yet
inseparable from it. The second (23 July) treated the free choice and moral test of the
angels, and the refusal expressed in the words \latin{Non ti servirò!}
(``I will not serve''). The third (30 July) turned to the
name: \latin{angelus} means messenger, and the Hebrew \latin{mal'ak} means
one ``delegated'', an ``ambassador''; the Pope drew the offices of the angels
from Ps 102[103], the Book of Tobias, and the princes of the nations in
Daniel~10.

The fourth audience (6 August) is the most doctrinal. It restated the
Lateran and Vatican definitions (DS~3002), described the angels as
immaterial, immortal, and personal, and named Raphael, Gabriel, and Michael
as the three whom Scripture names. It set out the nine choirs as Dionysius
arranges them and taught the ministry of the guardian angels, citing
Basil. The fifth (13 August) presented the fall of Satan by
the Lateran text, called him with the Gospel ``the prince of this world''
(\latin{principe di questo mondo}), and stated that possession is not
excluded. The sixth (20 August) closed the series: the power of Satan is
not infinite, Christ has conquered him, the Church exercises that victory
in exorcism, and the faithful pray \latin{Angele Dei} to their guardian.

\sectionguard
\subsection{The Catechism and the Compendium}

The \emph{Catechism of the Catholic Church} (1992; Latin typical edition
1997) gathers the whole of this doctrine for the universal Church.

\paragraph{The angels (CCC 328--336).} The Catechism opens with the
strongest statement of any modern act: ``The existence of the spiritual,
non-corporeal beings that Sacred Scripture usually calls `angels' is a
truth of faith. The witness of Scripture is as clear as the unanimity of
Tradition'' (CCC~328). It takes the distinction between office and nature
from Augustine (CCC~329; \emph{Enarr.\ in Ps.}\ 103, s.~1, 15). It teaches
that ``as purely spiritual creatures angels have intelligence and will:
they are personal and immortal creatures, surpassing in perfection all
visible creatures'' (CCC~330). It places them in relation to Christ:
``Christ is the centre of the angelic world'' (CCC~331); they were created
through and for him, and they serve the whole economy of salvation from
creation to the judgment (CCC~332--333). The whole life of the Church
``benefits from the mysterious and powerful help of angels'' (CCC~334); the
liturgy joins them in adoration and venerates certain of them (CCC~335).
Human life ``from infancy to death'' is surrounded by their care, and the
Catechism quotes Basil: ``Beside each believer stands an angel as protector
and shepherd leading him to life'' (CCC~336).

\paragraph{The fall of the angels (CCC 391--395).} The Catechism sees behind
the disobedience of our first parents ``a seductive voice, opposed to
God'', which Scripture and Tradition identify as a fallen angel called Satan
or the devil (CCC~391), and quotes DS~800. The fall ``consists in the free
choice of these created spirits, who radically and irrevocably rejected God
and his reign'' (CCC~392). ``It is the irrevocable character of their
choice, and not a defect in the infinite divine mercy, that makes the
angels' sin unforgivable'' (CCC~393), and the Catechism cites John
Damascene: there is no repentance for the angels after their fall, as there
is none for men after death. Scripture shows the devil as ``a murderer from
the beginning'' who tried to divert Jesus from his mission (CCC~394), yet
``the power of Satan is, nonetheless, not infinite'' (CCC~395): he is a
creature, and his action is permitted within a providence that brings good
out of it. The summary at CCC~414 gathers the teaching: ``Satan or the devil
and the other demons are fallen angels who have freely refused to serve God
and his plan. Their choice against God is definitive.''

\paragraph{The Compendium (2005).} The \emph{Compendium} approved by Benedict
XVI on 28 June 2005 condenses the same doctrine. The angels are ``purely
spiritual creatures, incorporeal, invisible, immortal, and personal beings
endowed with intelligence and will'' (\emph{Compendium} 60); they surround
Christ, serve his saving mission, and accompany human life with their
protection (\emph{Compendium} 61). The fall of the angels is the free and
irrevocable rejection of God by spirits created good, and the power of Satan
is not infinite (\emph{Compendium} 74--75).

\sectionguard
\subsection{The Church's Teaching by Degree}

The Church's teaching on the angels, each truth with its degree and the act
that teaches it:

\begin{threestable}{Grade}{Proposition}{Basis}
Defined & Spiritual creatures, the angelic, exist and were created by God
from nothing at the beginning of time & Lateran~IV, DS~800; Vatican~I,
DS~3002, 3025\\
Defined & The angels are not of God's substance and are not eternal &
Lateran~IV, DS~800; Vatican~I, DS~3025; cf.\ Braga~I, DS~455\\
Defined & The devil and the other demons were created good by nature and
became evil of themselves & Lateran~IV, DS~800; cf.\ Leo~I, DS~286; Braga~I,
DS~457; Innocent~III, DS~797\\
Defined & The punishment of the demons is perpetual & Lateran~IV, DS~801;
cf.\ Constantinople 543, DS~411\\
Church teaching & The existence of the angels is a truth of faith; the
existence of the demons is a dogmatic datum of the ordinary and universal
teaching & CCC~328; CDF study 1975; Paul~VI, 15 Nov.\ 1972\\
Church teaching & The angels are personal and immortal, with intellect and
will & \emph{Humani generis} 26, DS~3891; CCC~330; \emph{Compendium} 60\\
Church teaching & The fall was a free and irrevocable choice; the demons'
sin is unforgivable because of that character & CCC~392--393, 414\\
Church teaching & The demons create nothing and rule nothing by their own
authority; Satan's power is not infinite & Braga~I, DS~458, 462--463;
CCC~395\\
Church teaching & Each of the faithful has a guardian angel & CCC~336;
John Paul~II, 6 Aug.\ 1986\\
Disputed & Whether the angels were created together with the visible world
or before it & DS~800 \latin{simul}; \ST{I q.~61 a.~3}\\
Common teaching & Nine orders of angels, gathered from the names in
Scripture & Gregory, \emph{Hom.\ in Ev.}\ 34.7; \ST{I q.~108 a.~6};
John Paul~II, 6 Aug.\ 1986\\
Disputed & The ordering of the nine orders among themselves & Dionysius
and Gregory; \ST{I q.~108 a.~6}\\
\end{threestable}

The number of the angels, the object and instant of the first sin, and the
ordering of the orders are taught by the Fathers and Doctors, whose
teaching the following sections give. On the names of the angels the Church
has spoken: Scripture names Michael, Gabriel, and Raphael, and from the
Roman synod of 745 to the \emph{Directory on Popular Piety} of 2001 the
Church has forbidden the invention of others (Appendix~\ref{app:chronology};
Section~\ref{sec:devotion}; Appendix~\ref{app:names}).
